\section{System Frame}

The safety-net idea becomes clearer when viewed as a small set of design choices rather than as a single recognition model. In this article, the system frame has three choices: the safety net is based on speech, the first processing step runs on the drone side, and the result must be available in real time. These choices are motivated by the same requirement: nearby speech should be usable when a person observes that a drone may be entering a risky situation without depending on an external interface, network path, or delayed interpretation.

The first design choice is to focus on speech as the input to the safety net. A nearby person may shout an interruption, warning, or motion-related request while watching the drone operate. Rotor noise, microphone placement, motion, and distance all shape the audio before any software sees it. A speech-to-text system is useful when the interface needs readable language, for example for logging, display, or later human review. Prior drone speech studies show that transcript-based voice control can be possible, but they also show that rotor noise and platform conditions make word recognition fragile around drones~\cite{oneata19_interspeech,miesikowska2021uavnoise}. We therefore do not use speech-to-text as the default path. In short, noisy, or urgent interaction, recovering every word can be both difficult and unnecessary. For the safety net, the first question is simpler: is the nearby speech emergency-related, movement-related, or unsupported? A direct audio model can answer that question as a compact category, avoiding an extra transcript step before the safety net receives the speech-derived signal. This design is closer to limited-vocabulary speech recognition than to open-ended dictation, where the goal is to recognize a small set of useful speech categories under constrained conditions~\cite{warden2018speechcommands}.

The second design choice is to keep the speech-processing step on the drone side. Spoken interaction around a drone should not assume that a nearby person can open a phone application, pair with a controller, or reach a cloud speech service before communicating. These assumptions are especially weak for emergency-related speech, where the person who needs to speak may be a bystander rather than the trained operator. Local processing also keeps the microphone, rotor noise, and first interpretation in the same physical setting. This requirement aligns with prior work on microcontroller speech and embedded inference, which shows that small-footprint speech models and TinyML runtimes can support local audio decisions under memory and operator constraints~\cite{zhang2017helloedge,david2021tflm,lin2020mcunet,banbury2021mlperftiny}.

The third design choice is real-time speech processing. Speech that indicates a need to stop, move away, or hold position loses value if it is interpreted only after the scene has changed. The speech module should therefore operate at the pace of the interaction: audio is captured in short windows, processed locally, and converted into a safety-net category quickly enough for the drone system to use. This timing requirement is another reason to avoid making full transcription the default first step. The safety net benefits from a compact output that can be produced under rotor-noisy conditions before the speech becomes stale.

Together, these choices define how the safety net should work. The speech module is not expected to make the drone understand arbitrary language. It turns nearby speech into a local, timely, and limited category: emergency-related speech, movement-related speech, or unsupported input. The safety net then keeps that category separate from immediate drone behavior. This gives the drone system a practical starting point while preserving the boundary that speech should not become physical behavior by default.
